\section{Experiencia laboral}

\cventry{Primavera 2024-Actualmente}
{Human Behavior Lab, Texas A\&M University}
{Investigador posdoctoral}{}{}{}

\cventry{Invierno 2024, Otoño 2024}
{University of California Santa Barbara}
{Jefe de asistentes de enseñanza. Estadística para Economía}
{}{}{}

\cventry{2018-2019}
{Association of Supervisors of Banks of the Americas}
{Asistente de investigación}{}{}{}

\cventry{2018}
{Instituto Federal de Telecomunicaciones}
{Subdirector de Análisis de Competencia Económica}
{}{}{}

\cventry{2017-2018}
{Laboratorio Nacional de Políticas Públicas}
{Investigador Junior}
{}{}{}

\subsection{Otros Nombramientos}

\cventry{Verano 2023}{Dirección de Educación Financiera y Fomento Cultural, Banco de México}
{Becario de investigación en BANXICO}
{}{}
{Programa de Investigación de Verano (PIV)}

\cventry{2021}{Economics Department, UCSB}
{Asistente de investigación}{Professor Sevgi Yuksel}
{}
{Implementation and payment of an online experiment on the Base Rate Neglect effect.}

\cventry{2016-2017}{Division de Economía, CIDE}
{Asistente de investigación}
{Professor Alexander Elbittar}
{}
{Data analysis of  laboratory and field experiments}

\cventry{2012-2015}{Laboratorio de Conducta y Adaptación, UNAM}
{Asistente de investigación}
{Professor Arturo Bouzas Riaño}
{}
{Proyectos: En aprendizaje en ambientes Dinámicos y Conducta de Riesgo: Evaluación de Modelos de Elección}{}{}{}

%\cventry{2014-2015}{Project: En aprendizaje in Dynamic Environments}{Professor Arturo Bouzas Riaño}{}{}{Laboratory of Behavior and Adaptation}
%
%\cventry{2012-2013}{Risk Behavior: Election Models Evaluation}{Professor Arturo Bouzas Riaño}{}{}{Laboratory of Behavior and Adaptation}

%\cventry{2011-2012}{
%	Training in strategic skills for problem-solving and research evaluation based on heuristics}{Professor Carlos Santoyo Velasco}{}{}{
%%Proyecto:Investigación Supervisada sobre el Desarrollo y Contexto del Comportamiento Social\\* 
%Laboratory of Development and Context of Social Behavior}
%
%\cventry{2011-2012}{
%	Resolution Time of Arithmetic Sums 
%	 as function of Size and Distance Effects}{Dr. Julio Espinosa Rodríguez}{}{}{Laboratory of Acquisition and Development of Numeric Skills}
	 
	 